\chapter{Research Impact}
\label{chapter:impact}

% ============================================================================
% NOTE TO RYAN: Everything in \AI{...} below is Claude's orange rough draft,
% built from your notes + the verified impact frenzy (see impact_verified.md).
% NONE of it is final. Your original handwritten notes are preserved as
% %-comments directly above each AI block so you can rewrite in your own voice.
% Search "\AI{" to find every passage still needing your words.
% ============================================================================

\AI{The chapters of this thesis describe five methods for controlling diffusion models. This chapter steps back from the methods themselves to trace their downstream influence: the follow-up papers they enabled, the open-source tools and community ports built around them, and the broader research directions they helped open. The works are discussed in the order they appear in this thesis.}

\section{Peekaboo}

% --- RYAN'S ORIGINAL NOTE ---
% Peekaboo:
% this work was impactful because it showed people that diffusion foundation
% models could be used, even as a black box that predicts noise and without
% pulling apart the attention layers etc, to do more than just image generation.
% It inspired well cited papers such as the "your diffusion model is secretly a
% classifier" paper which was extremely well cited [ but I am unfamiliar so
% please take a look ]. It was also the first work to demonstrate RGBA
% text-to-image generation, albeit at low quality. The mechanism behind this
% paper's optimization was also directly used for diffusion illusions
% imeediateyly after the paper was written.
% [FACT-CHECK: "inspired" -> "cited by" (bib-only); "first RGBA" -> "first to
%  attempt"; "immediately after" -> ~13mo later. See impact_verified.md sec 1.]
\AI{Peekaboo demonstrated that a frozen text-to-image diffusion model could be repurposed for a discriminative task, zero-shot referring segmentation, while treating the model as a black box that predicts noise, without dissecting its attention layers or retraining any weights. Its lasting contribution was less the segmentation numbers than the demonstration that these foundation models encode far more than the image-generation objective they were trained on, and that this latent structure can be extracted through inference-time optimization alone. The same score-distillation optimization that drives Peekaboo's alpha-mask search was subsequently reused as the engine of Diffusion Illusions (\cref{chapter:diffillusions}), and the broader observation has been taken up across zero-shot segmentation, with works such as DiffCut~\cite{couairon2024diffcut} building directly on diffusion-feature segmentation. Peekaboo is also, to our knowledge, among the first works to attempt transparent (RGBA) generation from an RGB-only diffusion model, albeit at low quality, predating dedicated trained RGBA generators by over a year. It is cited by both diffusion-as-classifier papers that established the parallel observation on the discriminative side,} \citet{li2023diffusionclassifier} \AI{and} \citet{clark2023texttoimage}\AI{.}

\section{Diffusion Illusions}

% --- RYAN'S ORIGINAL NOTE ---
% Diffusion Illusions:
% It basically spawned its own subfield of research. It won the CVPR demo award,
% got featured as a SIGGRAPH paper, and was the focus of two youtube videos by
% popular science communication youtubers Steve Mould and Matt Parker. It won the
% best CVPR demo award for (2024? 2023?) [claude: please see the website it talks
% more about that]. it triggered other works such as Visual Anagrams (did it get
% an oral?), the one by Pascal (lookingglass) which got a CVPR oral, and the
% multi-scale diffusion thing that Daneil Geng also made that got featued on the
% back of every CVPR shirt in 2025 (i remember now it was calld the factorized
% diffusion paper). Illusion3D was similar to lookingglass, and was also in
% SIGGRAPH I think? [Please list other diffusion works here that were inspired by
% it too - do a research frenzy]. Diffusion Illusions essentially spawned a new
% subfield of research: using foundation models to create optical illusions.
% [FACT-CHECK: Demo Award = CVPR 2023 (not 2024). Visual Anagrams = CVPR'24 ORAL.
%  LookingGlass = CVPR'25 Oral (Pascal Chang, Disney). Factorized = CVPR 2024
%  shirt NOT 2025, ECCV'24 paper. Illusion3D = NOT SIGGRAPH (arXiv'24 + CVPR'25
%  Art Gallery). "spawned subfield" softened to "helped establish" (Visual
%  Anagrams is near-simultaneous). See impact_verified.md sec 2.]
\AI{Diffusion Illusions is among the most visible works in this thesis. It was published at SIGGRAPH~2024 and won a Demo Award at CVPR~2023, and its illusions were the subject of companion videos by the popular science-communication YouTubers Steve Mould and Matt Parker, both of which credit the work directly. Its broader significance is that it helped establish an active line of research on using foundation models to generate optical illusions, a problem that had previously relied on hand-crafted or GAN-based pipelines. It did so by showing that score-distillation optimization over a frozen diffusion model could automatically produce physically fabricable illusions across several classical illusion families.}

\AI{This direction has since been pursued by a steady stream of follow-up work. Visual Anagrams~\cite{geng2023visualanagrams}, a CVPR~2024 Oral, introduced a simpler zero-shot alternative for view-conditioned illusions; Factorized Diffusion~\cite{geng2024factorized} extended the idea to frequency-decomposed hybrid images; PTDiffusion~\cite{gao2025ptdiffusion} produced hidden-image illusions through training-free phase transfer; and the direction expanded from two dimensions into three with Illusion3D~\cite{feng2024illusion3d} and the CVPR~2025 Oral LookingGlass~\cite{chang2025lookingglass}, which generates anamorphic 3D illusions. Taken together, these works form a small but recognizable cluster of diffusion-based illusion generation that Diffusion Illusions helped seed.}

\section{MAGICK}

% --- RYAN'S ORIGINAL NOTE ---
% MAGICK: [was this work high impact? who cited it any why?] It was at the time,
% the largest general domain RGBA image dataset available to the public [ the PDF
% paper alone had more examples of RGBA images in it than the previous largest
% matting dataset did].
% [FACT-CHECK: "public" is load-bearing (LayerDiffuse 1M was never released).
%  6 papers materially use it: Trans-Adapter, TKG-DM, PrismLayers, PSDiffusion,
%  LayeringDiff, a matting work. ~20 cites, 6448 HF dl/mo. See impact_verified.md
%  sec 3.]
\AI{At the time of its release, MAGICK was the largest publicly available general-domain RGBA image dataset. It filled this gap deliberately, since the only larger general-domain RGBA collection of the era was never released to the public, and the largest released general-domain matting dataset contained only several hundred foreground objects. The contrast is stark enough that the paper itself notes that the matted images exhibited in its own figures already outnumber that prior dataset. Because no competing large-scale RGBA dataset has appeared since, MAGICK has become the de-facto reference dataset for transparent and layered image generation: it is used directly for training or evaluation by works such as Trans-Adapter~\cite{dai2025transadapter}, TKG-DM~\cite{tkgdm2025}, and PrismLayers~\cite{prismlayers2025}, and its images form a substantial fraction of the LayerBench benchmark introduced by the former. Its impact is therefore best measured not in citation count or media attention, since it was a standard conference poster, but in sustained practical adoption as research infrastructure, reflected in thousands of dataset downloads per month.}

\section{Go-with-the-Flow}

% --- RYAN'S ORIGINAL NOTE ---
% Go with the flow: My other most impactful work - a CVPR oral. It was at its
% time the best motion control model SOTA - by a long shot - and stayed that way
% for months. This is because of how simple it was to set up, letting us jump on
% the CogVideoX5B base model the week it was released. Go-with-the-Flow was easy
% to set up cause all I needed to do was swap out one noise for another, and fine
% tune. As a result, the open source community picked up the code base and within
% a few weeks of release people had built GWTF versions for AnimateDiff and
% HunyuanVideo. People also built GUI's and Comfy Nodes around it too (zeptaframe
% and kijai made some). It received over 1000 stars on Github, and became the
% adopted base model for many future works such as [ that physics paper charles
% hermann made about playing with physics and using GWTF to render it ] and [ I
% know there were like 5 other papers that did this too. find them and talk to me
% about them with links etc].
% In addition to the motion control aspect of this paper, many other works have
% continued to make use of, and in many cases improve, the warped noise itself
% [ please research frenzy here to see who did I know there are several ]. Some
% people extended it to 3d even [who?]
% Also, not sure where I should put this, but GWTF was also used in a demo in
% CVPR 2025? I think? With Yosun Chang and I that got an honorable mention.
% Please look into that too - research it and get back to me.
% [FACT-CHECK: 1081 stars (Eyeline-Labs). Ports by spacepxl. ZeptaFrame + kijai's
%  ComfyUI-VideoNoiseWarp. Charles Herrmann paper = WonderPlay (ICCV'25). Build-on
%  papers: EquiVDM, WonderPlay, HumANDiff, CameraNoise (MotionMatcher does NOT use
%  it - dropped). 3D = CameraNoise (ICML'26). Yosun Chang demo "AI3D Render" = Best
%  Demo Honorable Mention. ~107 cites. "SOTA for months" NOT verified - softened.
%  See impact_verified.md sec 4.]
\AI{Go-with-the-Flow, a CVPR~2025 Oral, is the most widely adopted method in this thesis. Its central design choice, replacing a video diffusion model's i.i.d.\ sampling noise with warped noise that carries motion structure and then fine-tuning the model to follow it, made motion control unusually easy to bolt onto a base model, requiring no architectural changes. This simplicity let the method ride the newest base models almost as they appeared, and it translated directly into open-source uptake. The official release accumulated over a thousand GitHub stars, and within weeks the community had ported the technique to other backbones, including AnimateDiff and HunyuanVideo, and wrapped it in interactive tools and ComfyUI nodes for noise warping. The method has since served as a building block for downstream systems well outside its original motion-transfer setting, such as WonderPlay~\cite{wonderplay2025}, which uses it as a renderer for physically simulated dynamic scenes.}

\AI{Beyond its use as a motion-control method, the warped-noise idea at the core of Go-with-the-Flow has itself become an object of study and extension. EquiVDM~\cite{liu2025equivdm} provides a theoretical account of why warped noise works, framing it through the lens of equivariance, while other works carry the construction into new regimes, for example reformulating the warp around 3D camera geometry to obtain camera-controllable generation~\cite{cameranoise2026}. Separately, the warped-noise sampler powered an interactive demo, presented jointly with Yosun Chang at CVPR~2025, that received a Best Demo Honorable Mention.}


\section{MotionV2V}

% --- RYAN'S ORIGINAL NOTE ---
% MotionV2V:
% This paper is rather new (CVPR 2026) - so it's not clear what impact it will
% have yet. But it was the first paper to demonstrate a new type of edit. Motion
% edits are a class of edit that, if we want to be able to manipulate any video,
% will eventuanly have to become mainstream. Did anyone cite it yet?
% [FACT-CHECK: YES already ~10 cites (e.g. TrajectoryMover). Bare "first" REFUTED
%  (MotionEditor'23, MotionFollower'24, ReVideo'24) -> scope to "first
%  video-to-video general-object motion editor". CVPR 2026 confirmed. See
%  impact_verified.md sec 5.]
\AI{MotionV2V is the most recent work in this thesis, accepted to CVPR~2026, so its impact is necessarily still emerging. Its contribution is to define and address motion editing as a distinct class of video edit: while earlier methods could edit appearance or follow a supplied trajectory~\cite{motioneditor,motionfollower,revideo2024}, MotionV2V is, to our knowledge, the first to perform true video-to-video motion editing on general objects. It modifies the trajectories already present in a real video while preserving its content, rather than regenerating the scene or restricting itself to human skeletons. We expect this class of edit to become increasingly important as video models mature toward general-purpose manipulation. Despite its recency, the work has already begun to be cited by follow-up research on trajectory-based video editing, which suggests the problem framing is resonating with the community.}


\chapter{Conclusion and Future Work}
\label{chapter:conclusion}

% --- RYAN'S ORIGINAL NOTE ---
% [restate intro here basically]
% Foundation models are very powerful for many tasks. But they aren't useful if
% we stick to just text to image or text to video. Our ability to harness and
% control these models is what makes them useful, and is the focus of my work.
\AI{Diffusion foundation models have made photorealistic image and video generation broadly accessible, but raw text-to-image and text-to-video generation expose only a narrow slice of what these models can do. The user supplies a short description and the model fills in everything else (composition, transparency, motion, semantics) from its prior, leaving a persistent gap between what the user intends and what the model produces. The thesis takes the position that closing this gap is what makes foundation models genuinely useful, and that the structure needed to close it already lives inside the frozen model, waiting to be controlled. The five works presented here (Peekaboo, Diffusion Illusions, MAGICK, Go-with-the-Flow, and MotionV2V) each extend control along a different axis: spatial semantics, multi-view geometry, transparency, temporal motion, and motion editing. Together they show that a frozen or lightly fine-tuned diffusion model can be steered toward tasks its training objective never directly anticipated.}

\section{Future Work}

% --- RYAN'S ORIGINAL NOTE ---
% Future work:
% In the future, I'll be working on video editing tasks - such as proof that an
% image-to-image general editor can be turned into a video-to-video general
% editor (like nanobanana but for video, look it up so you understand what that
% means claude).
% [Nano Banana = Gemini 2.5 Flash Image: a single general-purpose model that
%  takes an image + a text instruction and makes a targeted edit while preserving
%  the rest. The future-work goal: the same idea, but for video.]
\AI{The natural next step extends this line of work from controlling video \emph{generation} to controlling video \emph{editing} in full generality. Recent image models have shown that a single general-purpose editor can take an image together with a natural-language instruction and apply a targeted change (swapping an object, altering a background, adjusting lighting) while leaving the rest of the image untouched. No equivalent general editor yet exists for video: editing a clip still tends to mean regenerating it, which sacrifices the temporal coherence and fine detail of the original. MotionV2V (\cref{chapter:motionv2v}) is an early step in this direction for the specific case of motion, and the broader goal is to show that an instruction-driven image editor can be turned into an instruction-driven \emph{video} editor that handles any edit, whether appearance, motion, or both, while preserving everything the user did not ask to change. Achieving this would make video as directly and intuitively editable as text has made images, and would be a meaningful step toward general-purpose control over video foundation models.}
